% Body fragment. Requires amsmath, amssymb, amsthm, longtable, hyperref.
% theorem, proposition, definition, corollary, remark environments are assumed.
% Bibliography keys: riehl2016, depaiva1991, kellyPR.
\section{Category theory without a universe of arbitrary functions}
\label{sec:category-framework}

The relevant part of category theory is largely a calculus of typed composition,
explicit universal operations, and equations.  Its ordinary definitions and
standard theorems are documented in Riehl's text \cite{riehl2016}.  The issue for
iPRA is the representation of the data to which that calculus is applied.  A
finite table can be evaluated uniformly by a primitive recursive function.  An
arbitrary primitive recursive program cannot.  The distinction must be preserved
at every use of the words ``category of'', ``functor'', and ``natural''.

\subsection{Three distinct meanings of a categorical development}

A \emph{fixed PR presentation} consists of finitely many specified primitive
recursive predicates and operations, together with iPRA proofs of their laws.
For example, object codes are natural numbers satisfying $O(a)$; an arrow code
$f:a\to b$ satisfies $H(a,b,f)$; and parallel arrows are compared by
$E(a,b,f,g)$.  Identities and composition are total PR operations $i(a)$ and
$c(a,b,d,g,f)$, whose values on invalid arguments are immaterial.  On valid
arguments the axioms assert:
\begin{enumerate}
\item $E$ is an equivalence relation on each $H(a,b,-)$;
\item $i(a):a\to a$, and $g\circ f:a\to d$ whenever $f:a\to b$ and $g:b\to d$;
\item composition respects $E$ in both arguments;
\item $(h\circ g)\circ f=_{E}h\circ(g\circ f)$,
      $i(b)\circ f=_{E}f$, and $f\circ i(a)=_{E}f$.
\end{enumerate}
Every displayed law carries the relevant object and arrow guards.  For a
\emph{decidable presentation}, $O,H,E$ are PR predicates.  More generally $E$ may
be a first-order arithmetic formula with supplied iPRA proofs of these laws;
then equality need not be decidable, and algorithms that test it require a
separate decidable presentation or certificate interface.  One must not silently
use equality testing in the second version.

A functor is a specified pair of PR maps on object and arrow codes preserving
these laws.  A natural transformation is a specified PR map $a\mapsto\eta_a$
satisfying the typed naturality equations.  These are \emph{schemes over fixed
maps}: the symbols are fixed before the iPRA statement is formed.  We do not
quantify over arbitrary functions or insert a universal evaluator for their
codes.  If a category's laws are only hypotheses, the resulting theorem is a
conditional theorem of iPRA; an application becomes an unconditional theorem
only after those hypotheses have themselves been discharged.

A second, stronger uniformity is available for \emph{finite table presentations}.
Here the entire finite category, functor, and natural transformation are encoded
by natural numbers, their validation predicates are PR, and all table operations
are PR.  A theorem may then quantify over their codes in one iPRA formula.

A third meaning is the \emph{external category of PR constructions}.  Its objects
and arrows are individual specified PR definitions with supplied proofs.  It is
legitimate mathematical bookkeeping and a source of theorem schemes.  It is not
a first-order object that iPRA has been shown to contain.  In particular a
syntactic program code is not by itself a morphism equipped with a uniform total
PR evaluation operation.

\subsection{Explicit category tables and finite functor categories}

A finite category table has object set $[n]$, arrow set $[a]$, source and
target arrays of length $a$, an identity array of length $n$, and a composition
array with $a^2$ entries. A composition value is used only when the target of
the first arrow equals the source of the second. Check array dimensions and
entry bounds first; only afterwards check typing, unit laws, and associativity
on composable triples. The final validity predicate is a bounded Boolean
conjunction. Off-domain composition receives a fixed default; the reference
implementation uses a separate sentinel to make invalid lookups visible.

For categories $C,D$ with sizes $(n,a)$ and $(m,b)$, enumerate the
$m^n b^a$ possible pairs of object and arrow tables. Filter by preservation
of source, target, identity, and composition. The retained list is exactly
the list of functors, including the unique empty functor from the empty
category. Empty-table counts use $0^0=1$.

For each retained pair $F,G$, enumerate the $b^n$ possible component tables
and retain exactly those with
$\eta_x:Fx\to Gx$ and
\[
 \eta_y\circ F(f)=G(f)\circ\eta_x\qquad(f:x\to y).
\]
The resulting category $[C,D]$ has these functors as objects and the tagged
retained component tables as arrows. Its identities are
$(1_F)_x=1_{Fx}$ and its composition is
$(\mu\circ\eta)_x=\mu_x\circ\eta_x$.
If $\eta:F\to G$ and $\mu:G\to H$, then
\[
 (\mu_y\eta_y)F(f)=\mu_yG(f)\eta_x
                 =H(f)\mu_x\eta_x,
\]
using the two naturality equations and associativity in $D$. Thus composition
passes the retained-table test. Unit and associative laws hold at each
component by the corresponding laws in $D$. Canonical ranking in the retained
lists converts these pointwise equations into equalities of arrow codes.
The number of objects is at most $m^n b^a$, and the number of arrows is at
most $(m^n b^a)^2b^n$. Every scan is bounded by these explicit PR expressions.
The filtered-prefix invariant proves that each retained list is sound and
complete. This constructs a functor category uniformly for finite inputs,
rather than presupposing a universe of all arbitrary functors.

\subsection{Elementary categorical results are intuitionistic}

For completeness, several often-used deductions require no induction or
excluded middle.  If $g,h:b\to a$ are both two-sided inverses of $f:a\to b$, then
\[
 g=g\circ 1_b=g\circ(f\circ h)=(g\circ f)\circ h=1_a\circ h=h.
\]
If $r\circ s=1$, then $s$ is monic: from $s\circ u=s\circ v$, composition with
$r$ gives $u=v$.  Dually $r$ is epic.  If a monic $f:a\to b$ has a right inverse
$r:b\to a$, then
$f\circ(r\circ f)=(f\circ r)\circ f=f=f\circ1_a$; cancellation gives
$r\circ f=1_a$.  These equations also prove the dual statement.  Thus an
isomorphism is both monic and epic, and a monic retraction or epic section is an
isomorphism.  The converse ``monic and epic implies isomorphism'' is not valid in
an arbitrary category: the two-object category with only identities and one
arrow $a\to b$ is a finite counterexample.

Two terminal objects $t,u$ have unique arrows $p:t\to u$ and $q:u\to t$; terminal
uniqueness gives $qp=1_t$ and $pq=1_u$.  The same argument applies to two objects
satisfying the same universal property, by working with their universal
structure maps.  An explicit mediating operation is preferable when the result
will be used algorithmically: a formula asserting existence and uniqueness does
not automatically provide a PR selector.

\subsection{The category of finite sets has exponentials in iPRA}
\label{subsec:finset-exp}

Write $[n]=\{0,\ldots,n-1\}$.  An arrow $[n]\to[m]$ is the canonical table of its
$n$ entries, each below $m$.  Equivalently, for $m>0$ it is a base-$m$ code
$c<m^n$ and its value at $j<n$ is
\[
 \operatorname{ev}_{n,m}(c,j)
 =\left\lfloor c/m^j\right\rfloor\bmod m.
\]
All operations are totalized by zero outside their valid domains.  When $n=0$
there is one empty table, including when $m=0$; when $n>0,m=0$ there is no table.
Thus the number of tables is $m^n$ with the usual convention $0^0=1$.

Identity is the table $j\mapsto j$.  Composition is the table
$j\mapsto\operatorname{ev}(g,\operatorname{ev}(f,j))$.  Equality is entrywise
bounded equality, or literal equality of canonical table codes.  Closure and the
category laws follow by evaluating at an arbitrary $j<n$.  All searches and
finite sums used to encode these tables have displayed finite bounds.  Their
proofs use induction only on PR Boolean predicates about finite prefixes.

The terminal object is $[1]$, the initial object is $[0]$, the product is
$[n]\times[m]\cong[nm]$, and the coproduct is $[n+m]$.  Use
$(i,j)\mapsto im+j$ for the product, with the empty cases handled by the guards.
The pairing of two tables $f:[z]\to[n]$ and $g:[z]\to[m]$ has entries
$f(k)m+g(k)$.  Its projection equations and uniqueness are entrywise equations.
The coproduct mediator is the table obtained by concatenating its two input
tables; its two equations and uniqueness are equally direct.

\begin{theorem}[Finite exponentials]
The table category of finite sets is cartesian closed, with
\[
 [m]^{[n]}=[m^n].
\]
Evaluation and currying are primitive recursive uniformly in $m,n,z$ and the
input tables, and their two inverse laws are provable in iPRA.
\end{theorem}
\begin{proof}
Evaluation sends $(c,j)$ to the $j$th entry of table $c$.  If either $n=0$ or
$m^n=0$, its domain is empty; otherwise the digit formula applies.  Given a table
$h:[z]\times[n]\to[m]$, define
\[
 (\operatorname{curry}h)(k)
   =\sum_{j<n}h(k,j)m^j.
\]
For $m>0$ this number is below $m^n$, and its $j$th digit is $h(k,j)$.  The
finite-digit lemma is proved by induction on the table length, using quotient
and remainder at the last digit.  For $n=0$ the sum is zero and $m^0=1$, so the
same assertion is valid; if $m=0,n>0$, existence of $h$ forces $z=0$, making
currying the empty table.  Consequently
$\operatorname{ev}(\operatorname{curry}h(k),j)=h(k,j)$ on every valid input.
Conversely a canonical code $c<m^n$ is recovered by summing its digits, so
currying the evaluation of any $k:[z]\to[m^n]$ recovers $k$.  These two equations
are precisely the exponential universal property.  No unbounded family of
functions has been encoded: every represented function has a specified finite
domain.
\end{proof}

This example prevents an overly strong negative conclusion.  Cartesian closure
itself is compatible with iPRA.  The obstruction below concerns an exponential
containing every PR function on an \emph{unbounded} natural-number domain with
uniform PR evaluation.

\subsection{Finite limits, including equalizers and pullbacks}

A finite diagram is a finite table category $J$ and a table functor $D$ assigning
the finite set $[d_i]$ to each object $i$ and a function table $D(a)$ to each arrow
$a:i\to j$.  List all tuples
\[
 x=(x_i)_{i\in\operatorname{Ob}J},\qquad x_i<d_i,
\]
whose ambient number is $P=\prod_i d_i$, with the empty product equal to one.
Retain a tuple exactly when every arrow satisfies
\[
 D(a)(x_i)=x_j.                                      \tag{L}
\]
Let $L$ be the increasing list of retained tuple codes.  Its size is a PR number
$\ell\le P$.  Ranking within $L$ and selecting its $r$th element are bounded PR
operations; rank is only used on an element known to have been retained.
Coordinate selection gives a cone $p_i:[\ell]\to[d_i]$.

For any cone $q_i:[z]\to[d_i]$ and $t<z$, the tuple $(q_i(t))_i$ satisfies (L),
by the cone equations.  Its rank in $L$ therefore defines a table
$u:[z]\to[\ell]$.  Then $p_i u=q_i$ for every $i$.  If $v$ also has these
projection equations, the tuple selected by $v(t)$ has the same every
coordinate as the tuple selected by $u(t)$; finite tuple extensionality and
injectivity of ranking give $v(t)=u(t)$.  Thus $u=v$.

This proves the finite-limit universal property and supplies its mediator.
For the empty diagram the construction gives $[1]$.  A parallel pair produces
the list of points where its two tables agree, so equalizers require no
undecidable set formation.  A cospan gives pairs with equal images, so pullbacks
are obtained by the same bounded filter.  The proof is uniform in the encoded
diagram.  It never applies a compactness principle to infinitely many
constraints.

\subsection{Finite colimits by bounded equivalence closure}

For the same diagram form the disjoint union
\[
 S=\coprod_i[d_i],\qquad N=\sum_i d_i.
\]
Its elements are tagged pairs $(i,x)$.  Generate an undirected relation by
\[
 (i,x)\sim(j,D(a)(x))\quad(a:i\to j),
\]
and include all reflexive pairs.  Its equivalence closure can be computed by
finite reachability.  One explicit procedure is Warshall's recurrence on the
$N$-element set:
\[
 R_{k+1}(u,v)=R_k(u,v)\lor\bigl(R_k(u,k)\land R_k(k,v)\bigr),
 \qquad k<N.
\]
Here $R_0$ is the reflexive symmetric generating relation.  The elementary
invariant says that $R_k(u,v)$ holds exactly when there is a path from $u$ to $v$
whose internal vertices are below $k$.  The backward implication splits a path
at its visits to vertex $k$: a segment between successive visits can be removed,
so at most one visit is necessary.  The forward implication concatenates the
two witnessing paths.  Both directions admit bounded implementations: on $N$
vertices any walk can have repeated-vertex segments deleted until its length is
below $N$, and at most its initial length many deletion steps are needed.
For the invariant it suffices to use paths of length below $N$; a temporary
concatenation has length at most $2N$ and is shortened in at most $2N$ deletion
steps.  The path and its successive deletions are finite coded arrays; the
invariant is a bounded predicate.  Thus this familiar proof uses quantifier-free induction
after Booleanization, not induction on a formula with an unbounded existential
path quantifier.

It follows that $R_N$ is an equivalence relation containing the generators.
For each $u<N$ choose the least $v<N$ with $R_N(u,v)$.  This bounded minimum
exists because $R_N(u,u)$ holds.  List the distinct minima increasingly, and
identify their list with $[q]$.  The map taking $u$ to the rank of its minimum is
our quotient map.  In particular no abstract quotient type or choice axiom is
needed in the object theory.

A cocone $t_i:[d_i]\to[z]$ has equal values on every generating edge.  Equality
propagates along a finite path, so its value is constant on each $R_N$-class.
Define $v:[q]\to[z]$ by evaluating the cocone at the chosen least
representative.  Then $v\circ\iota_i=t_i$.  Any other such map agrees on every
class because every class has its displayed representative.  This proves the
colimit universal property with a PR mediator.  The empty diagram yields
$[0]$; a parallel pair yields a coequalizer, and a span yields a pushout.

The finite reachability bound is essential.  Replacing this finite relation by
an unbounded computably presented congruence would not justify a PR normal form
or an equality decision procedure.

\subsection{Yoneda as a first-order calculation}
\label{subsec:yoneda}

Let $C$ be a fixed PR-presented category.  A contravariant PR-presented family
$F$ consists of a PR carrier predicate $P(x,u)$, a specified equivalence relation
$u\simeq_x v$ on this carrier, and a PR action $F(f,u)$ for $f:x\to y$ and
$u\in F(y)$, giving an element of $F(x)$.  The action respects equivalence and
satisfies
\[
 F(1_x,u)\simeq_xu,\qquad
 F(g\circ f,u)\simeq_xF(f,F(g,u)).                    \tag{F}
\]
It also respects equality of arrow codes.  These are guarded arithmetic
formulas.  The representable family $h_A(x)=C(x,A)$ uses precomposition as its
action.

A specified natural transformation $\alpha:h_A\to F$ is a PR operation
$\alpha(x,f)$ taking $f:x\to A$ to an element of $F(x)$, respecting arrow
comparison, and satisfying
\[
 \alpha(y,f\circ k)\simeq_y F(k,\alpha(x,f))
 \quad(k:y\to x,\ f:x\to A).                         \tag{N}
\]

\begin{theorem}[Presentation-wise Yoneda]
For every specified $\alpha$ satisfying (N), put
$a=\alpha(A,1_A)$.  Then
\[
 \alpha(x,f)\simeq_x F(f,a)\qquad(f:x\to A).          \tag{Y}
\]
Conversely every $a\in F(A)$ determines a natural transformation by
$\alpha^a(x,f)=F(f,a)$.  Evaluation at $1_A$ and this construction are inverse up
to the specified equivalences.  All operations are PR for the fixed
presentation, and the argument is an intuitionistic iPRA derivation from its
laws.
\end{theorem}
\begin{proof}
Use (N) with $x=A$, $f=1_A$, and $k$ equal to the given arrow $f:x\to A$.
It gives
\[
 \alpha(x,1_A\circ f)\simeq_x F(f,\alpha(A,1_A)).
\]
The left identity law and the requirement that $\alpha$ respect arrow equality
replace the left-hand side by $\alpha(x,f)$, proving (Y).  Conversely (F) gives
\[
 \alpha^a(y,f\circ k)=F(f\circ k,a)
 \simeq_y F(k,F(f,a))=F(k,\alpha^a(x,f)),
\]
which is naturality.  Evaluation of $\alpha^a$ at $1_A$ gives
$F(1_A,a)\simeq_Aa$ by (F).  Evaluation followed by reconstruction sends
$\alpha$ to a pointwise equivalent transformation by (Y).  No function
extensionality axiom inside iPRA is required: pointwise equivalence is the
arithmetic definition of equality of the displayed transformations.
\end{proof}

For $F=h_B$, the element $a$ is an arrow $A\to B$ and the transformation is
$f\mapsto a\circ f$.  Thus every specified natural transformation between
representables is induced by a unique arrow, modulo hom equality.  This is the
full-faithfulness assertion in presentation-wise form.  If $C$ is a finite table
category and all $F(x)$ are finite tables, then transformations themselves are
finite tables, (N) is a bounded test, and the same proof yields a \emph{single
uniform theorem about their codes}.  Only in that finite setting have we
constructed the complete collection of transformations as an internal finite
set.  In the general setting we have the theorem scheme (Y), which is already
sufficient for many universal-property arguments.

\subsection{Finite Kan extensions, explicitly}

The finite constructions go beyond elementary category laws.  Let $C,D$ be
finite table categories, $K:C\to D$ a table functor, and $F:C\to\mathrm{FinSet}$
a finite-valued table functor.  For $d\in D$, the left Kan extension can be
constructed as a quotient of the finite set of triples
\[
 (c,k,x),\qquad k:Kc\to d,\quad x\in F(c),
\]
by the equivalence relation generated by
\[
 (c,k,x)\sim(c',k',F(a)x)
 \quad\text{if }a:c\to c'\text{ and }k'\circ K(a)=k.   \tag{K}
\]
Use the finite-colimit algorithm just proved.  For $v:d\to d'$, send a triple
to $(c,v\circ k,x)$.  It preserves every generator (K), by associativity, so it
descends to the chosen representative quotient.  Identity and composition are
checked on representatives.  This defines the table functor $\operatorname{Lan}_K
F:D\to\mathrm{FinSet}$ by PR operations.

Its universal transformation sends $x\in F(c)$ to the class $[c,1_{Kc},x]$.
It is natural because for $a:c\to c'$ the image of this class under
$\operatorname{Lan}_K F(K(a))$ is $[c,K(a),x]$, which equals
$[c',1_{Kc'},F(a)x]$ by generator (K).
For any finite-valued $G:D\to\mathrm{FinSet}$ and natural transformation
$\alpha:F\to GK$, define
\[
 \widehat\alpha_d[c,k,x]=G(k)(\alpha_c(x)).            \tag{KL}
\]
This is well-defined: if (K) holds, naturality of $\alpha$ gives
\[
 G(k')\alpha_{c'}F(a)x
   =G(k')G(K(a))\alpha_cx
   =G(k'K(a))\alpha_cx
   =G(k)\alpha_cx.
\]
Equality then propagates along the bounded paths used in the quotient.  Formula
(KL) is natural in $d$ because $G(vk)=G(v)G(k)$.  Conversely a natural
$\beta:\operatorname{Lan}_K F\to G$ restricts to
$\alpha_c(x)=\beta_{Kc}[c,1,x]$.  Naturality at $k:Kc\to d$ recovers (KL), while
(KL) at the identity recovers $\alpha$.  Thus restriction and (KL) are inverse,
proving the Kan universal property rather than merely constructing candidate
objects.

For the right Kan extension, $(\operatorname{Ran}_K F)(d)$ is the finite set of
tuples
\[
 (x_{c,k})_{(c,k)},\qquad k:d\to Kc,\quad x_{c,k}\in F(c),
\]
satisfying
\[
 F(a)x_{c,k}=x_{c',K(a)\circ k}\quad(a:c\to c').       \tag{KR}
\]
It is a bounded filtered product, so the finite-limit algorithm applies.  An
arrow $v:d\to d'$ sends a tuple $x$ to the tuple with $(c,k':d'\to Kc)$-coordinate
$x_{c,k'v}$.  Condition (KR) and the functor laws verify this reindexing.
Given $\alpha:GK\to F$, the corresponding transformation is
\[
 \widehat\alpha_d(y)_{c,k}=\alpha_c(G(k)y).
\]
Its coherence is naturality of $\alpha$.  The counit
$\varepsilon_c(x)=x_{c,1_{Kc}}$ is natural: for $a:c\to c'$, its two
naturality sides are $x_{c',K(a)}$ and $F(a)x_{c,1}$, equal by (KR).
Conversely a natural $\beta:G\to\operatorname{Ran}_K F$ gives
$\alpha_c=\varepsilon_c\beta_{Kc}$.  Naturality of $\beta$ at $k:d\to Kc$
shows
\[
 \alpha_c(G(k)y)
 =(\operatorname{Ran}_K F(k)(\beta_d y))_{c,1}
 =(\beta_d y)_{c,k}.
\]
Thus reconstructing $\beta$ recovers every coordinate; restricting the
constructed $\widehat\alpha$ at the identity recovers $\alpha$ by the identity
law of $G$.  Empty comma categories cause no difficulty: the corresponding product
is the singleton set and the corresponding coproduct quotient is empty.

For explicit bounds put $q=|\operatorname{Ob}C|\,|\operatorname{Mor}D|$
and $M=\max(1,\max_c|F(c)|)$, using $M=1$ when $C$ is empty.
The left construction has at most $qM$ triples; the right construction has
at most $M^q$ candidate tuples.  Every indexing set in these constructions is
therefore finite and explicitly bounded by the input table sizes.  Thus finite Kan extensions and
their adjunctions to restriction are available uniformly in iPRA.  This result
does not assert Kan extensions over arbitrary infinite categories.

\subsection{Adjunctions from unit and counit}

For fixed presented categories, let $L:C\to D$, $R:D\to C$, natural maps
$\eta:1_C\to RL$ and $\varepsilon:LR\to1_D$, and the triangle equations
\[
 \varepsilon_{Lx}\circ L(\eta_x)=1_{Lx},\qquad
 R(\varepsilon_y)\circ\eta_{Ry}=1_{Ry}.
\]
Define
\[
 \Phi(f)=R(f)\circ\eta_x\quad(f:Lx\to y),\qquad
 \Psi(g)=\varepsilon_y\circ L(g)\quad(g:x\to Ry).
\]
Then
\begin{align*}
 \Psi\Phi(f)
 &=\varepsilon_y\circ LR(f)\circ L(\eta_x)\\
 &=f\circ\varepsilon_{Lx}\circ L(\eta_x)=f,\\
 \Phi\Psi(g)
 &=R(\varepsilon_y)\circ RL(g)\circ\eta_x\\
 &=R(\varepsilon_y)\circ\eta_{Ry}\circ g=g.
\end{align*}
The second lines use naturality of counit and unit respectively; the last use
the triangles.  For $a:x'\to x$, $b:y\to y'$, the identity
\[
 \Phi(b\circ f\circ L(a))=R(b)\circ\Phi(f)\circ a
\]
follows by expanding $\Phi$, using functoriality, and naturality of $\eta$.
Thus these are natural inverse hom correspondences, giving the adjunction.
This proof is a finite equational derivation in iPRA for the fixed signatures;
it assumes neither hom-set enumeration nor a function-space object.

Conversely, \emph{supplied} PR natural inverse hom correspondences give
$\eta_x=\Phi(1_{Lx})$ and $\varepsilon_y=\Psi(1_{Ry})$.  Naturality and the inverse
laws recover the formulas above, and hence the triangle equations.  A bare
assertion that some correspondence exists is insufficient for this
algorithmic conclusion unless its maps or bounded selectors are provided.

\subsection{Equivalence requires explicit splittings}

Suppose $F:C\to D$ is faithful and has a supplied PR lifting operation
$\lambda_{a,b}:D(Fa,Fb)\to C(a,b)$ with
$F(\lambda(h))=h$.  Suppose also that for each object $d$ we are supplied PR data
\[
 c_d\in C,\qquad s_d:Fc_d\to d,\qquad t_d:d\to Fc_d,
 \quad s_dt_d=1_d,\quad t_ds_d=1_{Fc_d}.
\]
Define
\[
 G(d)=c_d,\qquad
 G(f:d\to e)=\lambda_{c_d,c_e}(t_e\circ f\circ s_d).
\]
Applying $F$ shows that $G$ preserves identities and composition, because
$s_et_e=1_e$.  Faithfulness then proves the original equalities in $C$.
The arrows $s_d$ are the components of a natural isomorphism $FG\cong1_D$:
\[
 s_e\circ F(Gf)=s_et_efs_d=fs_d.
\]
The components
\[
 \eta_c=\lambda_{c,c_{Fc}}(t_{Fc}),\qquad
 \eta_c^{-1}=\lambda_{c_{Fc},c}(s_{Fc})
\]
are inverse by faithfulness.  Naturality also follows by applying $F$ and
canceling $s_{Fc}t_{Fc}$ in the resulting equation.  Hence $1_C\cong GF$,
providing an explicit equivalence.

This is a constructive \emph{split full, faithful, essentially-surjective}
criterion.  In finite table categories all missing data can be chosen by
bounded search over the finite object and arrow tables, so the familiar
full--faithful--essentially-surjective criterion is a uniform iPRA theorem
there.  In an infinite presentation, the formulas expressing fullness and
essential surjectivity alone do not furnish those PR selectors.  The standard
set-theoretic proof makes choices at precisely this point.  Replacing the
missing choices by supplied data or proved bounded searches states the actual
proof obligation.

\subsection{The exact diagonal obstruction}

\begin{proposition}[No universal PR evaluation on natural codes]
There is no total PR function $U(e,n)$ such that every unary PR function $f$ is
represented by some $e$ with $U(e,n)=f(n)$ for every $n$.
\end{proposition}
\begin{proof}
If $U$ were PR, then $d(n)=U(n,n)+1$ would be PR.  Let $e$ represent $d$.
Evaluating at $n=e$ gives $U(e,e)=U(e,e)+1$, impossible.  The contradiction uses
only successor's elementary arithmetic properties.  The representation
assertion ranges over all PR functions externally; for any fixed proposed
$U$, the contradictory instance is the one fixed PR definition $d$.
\end{proof}

The standard indexing-and-diagonalization mechanism is discussed in Kelly's
lecture notes \cite{kellyPR}; the displayed argument specifies the exact
interface ruled out here.  More generally suppose a proposed function object
has natural-number codes $E\subseteq\mathbb N$, evaluation is the restriction of
a total PR map $U$ on all natural inputs, and every unary PR function has a code
in $E$.  The same argument chooses $e\in E$ representing $d$, so restricting
the valid codes does not help.

Consequently a category of PR-presented subsets of the naturals with morphisms
given by restrictions of total PR maps cannot have an exponential
$\mathbb N^{\mathbb N}$ representing all such maps and evaluated by a morphism
of that same kind.  This is a limitation of that particular representation,
not a claim that iPRA proves no cartesian closed categories exist.  Finite sets
already provide the counterexample to that stronger claim.

There are three sound alternatives.  First, fix each PR definition before
forming its iPRA statement.  Second, use a primitive recursive bounded
interpreter $U(e,n,t)$, with a supplied execution bound or a checked finite
trace.  Third, use a uniform PR evaluator for a genuinely bounded fragment,
with the fragment's bound part of the interface.  In the latter two cases a
proof that every intended input is covered is an additional obligation.  A
finite proof checker does not supply a uniform soundness theorem for all its
accepted proofs, nor a uniform termination bound for all PR program codes.
